% ============================================================================
% TECHNICAL APPENDIX 14A: Fairness Metrics, Bias Auditing, and Algorithmic Accountability
% Source: /chapters/ch14/Ch14A_final.md
% ============================================================================

\chapter*{Technical Appendix 14A: Fairness Metrics, Bias Auditing, and Algorithmic Accountability Tools}
\addcontentsline{toc}{chapter}{Technical Appendix 14A: Fairness Metrics and Algorithmic Accountability}

\begin{chapterquote}{Mathematical foundations and Python implementations for building and auditing fair AI systems}
Fairness in machine learning is not a single property---it is a family of distinct mathematical requirements, some mutually incompatible. This appendix formalizes the major definitions, proves the impossibility theorem, and provides implementation tools for bias auditing and mitigation.
\end{chapterquote}

% ============================================================================
\section{Formal Fairness Definitions}
\label{sec:a14-fairness-definitions}
% ============================================================================

Let:
\begin{itemize}
    \item $\hat{Y}$ = model prediction (binary: 0 or 1)
    \item $Y$ = true outcome
    \item $A$ = sensitive attribute (e.g., race, gender, age group)
    \item $X$ = all other features
\end{itemize}

\paragraph{Demographic Parity (Statistical Parity).}
\[
P(\hat{Y} = 1 \mid A = 0) = P(\hat{Y} = 1 \mid A = 1)
\]
The probability of a positive prediction is equal across groups. A hiring algorithm satisfies demographic parity if it recommends candidates at equal rates regardless of protected attribute.

\paragraph{Equalized Odds \autocite{hardt2016equality}.}
\[
P(\hat{Y} = 1 \mid Y = y, A = 0) = P(\hat{Y} = 1 \mid Y = y, A = 1) \quad \text{for } y \in \{0,1\}
\]
Both the true positive rate (TPR) and false positive rate (FPR) are equal across groups. A medical screening tool satisfies equalized odds if it correctly identifies sick patients---and incorrectly flags healthy ones---at equal rates regardless of demographic group.

\paragraph{Equal Opportunity} (weaker form of equalized odds): only the TPR constraint:
\[
P(\hat{Y} = 1 \mid Y = 1, A = 0) = P(\hat{Y} = 1 \mid Y = 1, A = 1)
\]

\paragraph{Calibration (Predictive Rate Parity).}
\[
P(Y = 1 \mid \hat{Y} = 1, A = 0) = P(Y = 1 \mid \hat{Y} = 1, A = 1)
\]
Among all individuals predicted positive, the fraction who are actually positive is equal across groups.

% ============================================================================
\section{The Fairness Impossibility Theorem}
\label{sec:a14-impossibility}
% ============================================================================

\textcite{chouldechova2017fair} proved that calibration and equal FPR cannot simultaneously hold when base rates differ across groups.

\paragraph{Setup.} Let $b_A = P(Y = 1 \mid A)$ denote the base rate of the outcome for group $A$.

\paragraph{Theorem \autocite{chouldechova2017fair}.} If $b_0 \neq b_1$ (base rates differ), then no score $s$ can simultaneously satisfy:
\begin{enumerate}
    \item \textbf{Calibration:} $P(Y=1 \mid s, A) = s$ for all groups
    \item \textbf{Equal FPR:} $P(\hat{Y}=1 \mid Y=0, A=0) = P(\hat{Y}=1 \mid Y=0, A=1)$
    \item \textbf{Equal FNR:} $P(\hat{Y}=0 \mid Y=1, A=0) = P(\hat{Y}=0 \mid Y=1, A=1)$
\end{enumerate}

\paragraph{Intuitive proof.} For a calibrated model at threshold $\tau$:
\[
\text{FPR}_A = \frac{(1-\text{PPV}_A) \cdot \text{PR}_A}{1 - b_A}
\]
where $\text{PPV}_A$ is the positive predictive value and $\text{PR}_A$ is the prediction rate for group $A$. If $b_0 \neq b_1$ and PPV is equalized (calibration holds), then FPR equality requires $\text{PR}_A$ to adjust in a way that inevitably violates FNR equality. All three cannot hold simultaneously.

\begin{lstlisting}[style=python, caption={Fairness impossibility demonstration}]
import numpy as np
from sklearn.metrics import confusion_matrix

def compute_fairness_metrics(y_true, y_pred, groups):
    results = {}
    for group_val in [0, 1]:
        mask = groups == group_val
        yt, yp = y_true[mask], y_pred[mask]
        tn, fp, fn, tp = confusion_matrix(yt, yp, labels=[0, 1]).ravel()
        n = len(yt)
        results[f'group_{group_val}'] = {
            'TPR':       tp / (tp + fn) if (tp + fn) > 0 else np.nan,
            'FPR':       fp / (fp + tn) if (fp + tn) > 0 else np.nan,
            'PPV':       tp / (tp + fp) if (tp + fp) > 0 else np.nan,
            'pred_rate': (tp + fp) / n,
            'base_rate': (tp + fn) / n,
        }
    return results

def demonstrate_impossibility(base_rate_0=0.2, base_rate_1=0.4, n=2000):
    np.random.seed(42)
    y0 = np.random.binomial(1, base_rate_0, n // 2)
    y1 = np.random.binomial(1, base_rate_1, n // 2)
    y  = np.concatenate([y0, y1])
    groups = np.concatenate([np.zeros(n // 2), np.ones(n // 2)])

    scores = np.where(
        groups == 0,
        np.random.beta(2, 8, n),
        np.random.beta(4, 6, n)
    )
    y_pred  = (scores > 0.3).astype(int)
    metrics = compute_fairness_metrics(y, y_pred, groups)

    for g in ['group_0', 'group_1']:
        m = metrics[g]
        print(f"{g}: TPR={m['TPR']:.3f}, FPR={m['FPR']:.3f}, "
              f"PPV={m['PPV']:.3f}, PredRate={m['pred_rate']:.3f}")

    fpr_gap = abs(metrics['group_0']['FPR'] - metrics['group_1']['FPR'])
    ppv_gap = abs(metrics['group_0']['PPV'] - metrics['group_1']['PPV'])
    print(f"\nFPR gap: {fpr_gap:.3f} | PPV gap (calibration): {ppv_gap:.3f}")
    print("Note: minimising PPV gap (calibration) cannot prevent FPR gap "
          "when base rates differ.")
\end{lstlisting}

% ============================================================================
\section{Bias Detection with Fairlearn}
\label{sec:a14-fairlearn}
% ============================================================================

Microsoft's Fairlearn library provides audit dashboards and mitigation algorithms.

\begin{lstlisting}[style=python, caption={Full fairness audit with Fairlearn}]
from fairlearn.metrics import (
    MetricFrame, demographic_parity_difference,
    equalized_odds_difference, selection_rate,
)
from sklearn.metrics import accuracy_score, precision_score, recall_score

def full_fairness_audit(y_true, y_pred, sensitive_features, feature_name):
    mf = MetricFrame(
        metrics={
            'accuracy':       accuracy_score,
            'precision':      precision_score,
            'recall':         recall_score,
            'selection_rate': selection_rate,
        },
        y_true=y_true,
        y_pred=y_pred,
        sensitive_features=sensitive_features,
    )
    dp_diff = demographic_parity_difference(
        y_true, y_pred, sensitive_features=sensitive_features
    )
    eo_diff = equalized_odds_difference(
        y_true, y_pred, sensitive_features=sensitive_features
    )

    print(f"Fairness Audit: {feature_name}")
    print(mf.by_group.to_string())
    print(f"Demographic parity diff:  {dp_diff:.4f}  (0 = fair)")
    print(f"Equalized odds diff:      {eo_diff:.4f}  (0 = fair)")

    if abs(dp_diff) > 0.1:
        print(f"WARNING: Demographic parity gap ({dp_diff:.3f}) exceeds 0.1")
    if abs(eo_diff) > 0.1:
        print(f"WARNING: Equalized odds gap ({eo_diff:.3f}) exceeds 0.1")

    return {'metric_frame': mf, 'dp_diff': dp_diff, 'eo_diff': eo_diff}
\end{lstlisting}

% ============================================================================
\section{Bias Mitigation Techniques}
\label{sec:a14-mitigation}
% ============================================================================

Three categories of mitigation correspond to three pipeline stages: \textbf{pre-processing} (modify training data before training), \textbf{in-processing} (add fairness constraints to the learning objective), and \textbf{post-processing} (adjust predictions after training).

\begin{lstlisting}[style=python, caption={In-processing, post-processing, and pre-processing mitigation}]
from fairlearn.reductions import ExponentiatedGradient, DemographicParity
from fairlearn.postprocessing import ThresholdOptimizer
from sklearn.linear_model import LogisticRegression
import numpy as np

# In-processing: Exponentiated Gradient
def train_fair_inprocessing(X_train, y_train, sensitive_train):
    mitigator = ExponentiatedGradient(
        estimator=LogisticRegression(max_iter=500),
        constraints=DemographicParity(),
        eps=0.01,
    )
    mitigator.fit(X_train, y_train, sensitive_features=sensitive_train)
    return mitigator

# Post-processing: Threshold Optimizer
def train_fair_postprocessing(estimator, X_train, y_train, sensitive_train):
    postprocess_est = ThresholdOptimizer(
        estimator=estimator,
        constraints='equalized_odds',
        predict_method='predict_proba',
        objective='accuracy_score',
    )
    postprocess_est.fit(X_train, y_train, sensitive_features=sensitive_train)
    return postprocess_est

# Pre-processing: Reweighing (removes group-label correlation by instance weights)
def compute_reweighing_weights(y_train, sensitive_train):
    n = len(y_train)
    weights = np.ones(n)
    for g in np.unique(sensitive_train):
        for l in np.unique(y_train):
            mask  = (sensitive_train == g) & (y_train == l)
            n_gl  = mask.sum()
            if n_gl > 0:
                expected  = (sensitive_train == g).sum() / n \
                          * (y_train == l).sum() / n * n
                weights[mask] = expected / n_gl
    return weights
\end{lstlisting}

% ============================================================================
\section{Disparate Impact Ratio (AIF360)}
\label{sec:a14-dir}
% ============================================================================

The disparate impact ratio (DIR) is the standard metric used in US employment discrimination law (the ``4/5ths rule''):
\[
\text{DIR} = \frac{P(\hat{Y}=1 \mid A=\text{minority})}{P(\hat{Y}=1 \mid A=\text{majority})}
\]
A DIR below 0.8 is a legal threshold for disparate impact in employment contexts.

\begin{lstlisting}[style=python, caption={Disparate impact ratio via AIF360}]
from aif360.datasets import BinaryLabelDataset
from aif360.metrics import BinaryLabelDatasetMetric

def compute_disparate_impact(df, label_col, sensitive_col,
                              privileged_val, unprivileged_val):
    privileged_groups   = [{sensitive_col: privileged_val}]
    unprivileged_groups = [{sensitive_col: unprivileged_val}]

    dataset = BinaryLabelDataset(
        df=df[[sensitive_col, label_col]],
        label_names=[label_col],
        protected_attribute_names=[sensitive_col],
    )
    metric = BinaryLabelDatasetMetric(
        dataset,
        unprivileged_groups=unprivileged_groups,
        privileged_groups=privileged_groups,
    )
    dir_val   = metric.disparate_impact()
    mean_diff = metric.mean_difference()

    print(f"Disparate Impact Ratio: {dir_val:.4f}")
    print(f"  (Legal threshold: >= 0.80; below = potential disparate impact)")
    print(f"Mean Difference (stat parity): {mean_diff:.4f}")
    if dir_val < 0.80:
        print("ALERT: Disparate impact below 4/5ths rule threshold.")
    return {'disparate_impact': dir_val, 'mean_difference': mean_diff}
\end{lstlisting}

% ============================================================================
\section{Counterfactual Fairness}
\label{sec:a14-counterfactual}
% ============================================================================

Counterfactual fairness \autocite{kusner2017counterfactual} asks: would the prediction change if the individual belonged to a different demographic group, holding all causally independent features constant?
\[
P(\hat{Y}_{A \leftarrow a}(U) = y \mid X=x, A=a) = P(\hat{Y}_{A \leftarrow a'}(U) = y \mid X=x, A=a)
\]
This requires a causal model of the data-generating process. The testing principle can be approximated by flipping the sensitive attribute and measuring prediction change.

\begin{lstlisting}[style=python, caption={Approximate counterfactual fairness test}]
def counterfactual_fairness_test(model, X_test, sensitive_col,
                                  group_a_val, group_b_val,
                                  non_sensitive_cols):
    X_a         = X_test[X_test[sensitive_col] == group_a_val].copy()
    X_a_counter = X_a.copy()
    X_a_counter[sensitive_col] = group_b_val

    preds_original = model.predict(X_a[non_sensitive_cols])
    preds_counter  = model.predict(X_a_counter[non_sensitive_cols])

    changed = (preds_original != preds_counter).mean()
    print(f"Counterfactual flip rate: {changed:.4f}")
    print(f"  (0 = counterfactually fair; >0 = predictions changed by group flip)")
    return changed
\end{lstlisting}

% ============================================================================
\section{Feedback Loop Detection}
\label{sec:a14-feedback-loops}
% ============================================================================

Detecting whether a deployed model is amplifying disparities over time requires monitoring group-level prediction rates across retraining cycles.

\begin{lstlisting}[style=python, caption={Feedback loop amplification test}]
import numpy as np
from scipy import stats

def detect_feedback_amplification(predictions_history, group_labels,
                                   n_cycles=5):
    disparity_over_time = []
    for cycle in range(min(n_cycles, len(predictions_history))):
        preds  = np.array(predictions_history[cycle])
        groups = np.array(group_labels[cycle])
        rates  = {g: preds[groups == g].mean() for g in np.unique(groups)}
        disparity_over_time.append(max(rates.values()) - min(rates.values()))

    cycles = np.arange(len(disparity_over_time))
    slope, _, r_val, p_val, _ = stats.linregress(cycles, disparity_over_time)

    print(f"Disparities over cycles: {[f'{d:.4f}' for d in disparity_over_time]}")
    print(f"Trend slope: {slope:.4f}/cycle, R²={r_val**2:.3f}, p={p_val:.4f}")
    if slope > 0 and p_val < 0.05:
        print("WARNING: Statistically significant disparity amplification.")
    elif slope > 0:
        print("NOTE: Upward trend in disparity (not yet statistically significant).")
    else:
        print("No evidence of feedback amplification detected.")

    return {'slope': slope, 'p_value': p_val,
            'amplification_detected': slope > 0 and p_val < 0.05}
\end{lstlisting}

% ============================================================================
\section{Pre-Deployment Audit Checklist}
\label{sec:a14-checklist}
% ============================================================================

A structured audit checklist derived from NIST AI RMF (2023) and EU AI Act requirements:

\begin{center}
\begin{tabular}{@{}lp{4cm}p{3cm}p{3.5cm}@{}}
\toprule
\textbf{Stage} & \textbf{Check} & \textbf{Tool} & \textbf{Threshold} \\
\midrule
Data       & Demographic representation in training set & EDA + sampling stats  & No group $<$ 5\% of $N$ \\
Data       & Label distribution by group                & MetricFrame           & Parity diff $<$ 0.10 \\
Model      & Demographic parity at deployment threshold & Fairlearn             & Diff $<$ 0.10 \\
Model      & Equalized odds                             & Fairlearn             & Diff $<$ 0.10 \\
Model      & Calibration by group                       & ECE per group         & Per-group ECE $<$ 0.05 \\
Model      & Disparate impact ratio                     & AIF360                & DIR $\geq$ 0.80 \\
Deployment & Slice-based monitoring (monthly)           & SliceFinder           & Alert at 5\% degradation \\
Deployment & Feedback loop correlation test             & Custom                & Slope $p < 0.05 \to$ halt \\
Process    & Reviewer demographic diversity             & Audit log             & Document + review \\
Process    & Contestability mechanism functional        & Process audit         & 100\% cases have pathway \\
\bottomrule
\end{tabular}
\end{center}

% ============================================================================
\section{Further Reading}
\label{sec:a14-further-reading}
% ============================================================================

\begin{itemize}
    \item \textcite{chouldechova2017fair} --- Fair prediction with disparate impact. \textit{Big Data, 5}(2).
    \item \textcite{hardt2016equality} --- Equality of opportunity in supervised learning. \textit{NeurIPS 29}.
    \item \textcite{kusner2017counterfactual} --- Counterfactual fairness. \textit{NeurIPS 30}.
    \item Bird, S., et al.\ (2020). Fairlearn: A toolkit for assessing and improving fairness in AI. \textit{Microsoft Research}.
\end{itemize}
